
\documentclass[border=4pt]{standalone}
\usepackage[utf8]{inputenc}
\usepackage[T1]{fontenc}
\usepackage{amsmath,amssymb,xcolor}
\usepackage{tikz}
\usetikzlibrary{arrows.meta,shapes.misc}
\renewcommand{\familydefault}{\sfdefault}
\definecolor{vbleu}{HTML}{1C5FD6}
\definecolor{orange}{HTML}{E8590C}
\definecolor{vert}{HTML}{2F9E44}
\definecolor{gris}{HTML}{495057}
\definecolor{violet}{HTML}{7048E8}
\tikzset{
  vec/.style={->,very thick,>={Stealth[length=2.4mm]}},
  ref/.style={orange,thick,dashed},
  par/.style={vert,thick,dashed},
  sec/.style={gris,thick,dashed},
  ptcross/.style={cross out,draw,thick,minimum size=1.6mm,inner sep=0pt},
}
\usepackage{array}
\setlength{\parindent}{0pt}\setlength{\parskip}{5pt}
\newcommand{\V}[1]{\overrightarrow{\mathrm{#1}}}
\begin{document}

\begin{minipage}{8.5cm}\large \colorbox{orange!8}{\begin{minipage}{8.2cm}\vspace{4pt}{\color{orange}\large Point méthode}\par Deux vecteurs sont colinéaires s'ils ont la \textit{même direction}.\par On trace la droite (AB), puis les droites (CD), (EF), (GH) et (AK).\par Un vecteur est colinéaire à $\V{AB}$ lorsque la droite qui le porte est parallèle à la droite (AB), ou confondue avec elle.\vspace{4pt}\end{minipage}}\end{minipage}
\end{document}
